\documentclass[10pt,a4paper]{amsart}
\usepackage[margin=19mm]{geometry}
\usepackage{amsmath,amssymb,mathtools}
\usepackage[scr=rsfs,cal=euler]{mathalfa}
\usepackage{xfrac}
\usepackage[unicode,hidelinks,bookmarks=false]{hyperref}
\newcommand{\ie}{i.e.\ }
\newcommand{\cf}{cf.\ }
\newcommand{\resp}{respectively\ }
\newcommand{\Subsection}{\subsection}
\providecommand{\nobreakdash}{\nobreak\hbox{-}}
\setlength{\emergencystretch}{2em}
\multlinegap0pt
\usepackage{amssymb}
\usepackage{mathtools}
\usepackage{tikz}
\newtheorem{thm}{Theorem}
\newcommand{\Cay}[2]{\mathrm{Cay}(#1,#2)}		%Cayley graph
\title{Eternal domination in Cayley graphs}
\author{MacKenzie Carr, Nancy E. Clarke, Gary MacGillivray, Joy Morris}
\date{Source: arXiv:2607.04024v1 (CC BY 4.0)}
\begin{document}
\maketitle

\noindent\emph{Source and licence.} Eternal domination in Cayley graphs. Version arXiv:2607.04024v1 (CC BY 4.0), distributed by arXiv under the Creative Commons Attribution 4.0 International licence, http://creativecommons.org/licenses/by/4.0/, as recorded in the arXiv licence metadata for this exact version and shown on the abstract page https://arxiv.org/abs/2607.04024v1. Full licence notice: \texttt{licenses/arxiv-2607.04024-CC-BY-4.0.txt}.\par\smallskip

\noindent\textit{Editorial scope.} Counted unit: the article's thm abelian (source0.tex lines 141--143). The article's complete original proof of it is delivered with it (source0.tex lines 145--165, one \texttt{proof} environment of 21 lines). No mathematics is written, repaired or re-derived: the statement and the proof are the article's own lines, in the article's own notation, and no retained line is substituted or re-flowed.  No further source-local context is needed: every cross-reference of every retained line resolves inside the two delivered spans.  Nothing was omitted from any retained span, so the package carries no omission record.  No retained line contains a citation, so the wrapper carries no bibliography and no bibliography entry.  No retained line contains a figure environment, an image-inclusion command or drawing code, so no image is delivered and the package declares no asset dependency.  Editorial formatting only, in addition: an AMS \texttt{amsart} preamble that restates the article's own declarations and macros used by the retained lines, namely the environments \texttt{thm} on separate counters, together with the standard packages those lines need.  The retained environments print in this document as Theorem 1 in order of appearance, whereas the article prints the same items with the numbering of its own full document.  The complete native source of the article as distributed in the arXiv e-print for this exact version is bundled unchanged as \texttt{sources/204441/source0.tex}.
\begin{thm}\label{abelian}
    Let $D$ be a dominating set for the Cayley graph $\Cay{G}{S}$, and let $s \in S$ be such that $s$ is either centralised or inverted by every element of $G$. Then for any vertex $v$ in $Ds$ there is a dominating set $D'$ containing $v$ such that guards can be reconfigured from $D$ to $D'$.
\end{thm}

\begin{proof}
    Observe that since every element of $G$ either centralises or inverts $s$, we can partition $G$ into $G_1=C_G(s)$ and $G_2=G \setminus G_1$. Furthermore, either $G=G_1$ or $G_1$ has index $2$ in $G$. Similarly, we can partition $D$ into $D_1=D\cap G_1$ and $D_2=D\cap G_2$, and $S$ into $S_1=S\cap G_1$ and $S_2=S\cap G_2$.

    Let $v=ds$ where $d \in D_i$, with $i \in \{1,2\}$. Then for every $d' \in D_i$, we move the defender from $d'$ to $d's$, while for every $d' \in D_{3-i}$, we move the defender from $d'$ to $d's^{-1}$ (we can do this since $S$ is inverse-closed). We claim that the vertices now occupied by the defenders form a dominating set $D'$ for $\Cay{G}{S}$.

    Observe that $D$ being a dominating set for $\Cay{G}{S}$ is equivalent to saying $G = DS$. (Here we are using the assumption that $1 \in S$.) We can refine this using our partitions to the following equivalent statement: $G_1=D_1S_1 \cup D_2S_2$ and $G_2=D_1S_2 \cup D_2S_1$.

    Although the proofs are very similar, we analyze the situations $i=1$ and $i=2$ separately. 

    Suppose first $i=1$. Then after moving, the defenders occupy the vertices $D_1s \cup D_2s^{-1}$, and $$(D_1s \cup D_2 s^{-1})S  = D_1sS_1 \cup D_1sS_2 \cup D_2s^{-1}S_1 \cup D_2s^{-1}S_2 $$
    $$ = D_1S_1s \cup D_1S_2s^{-1} \cup D_2S_1s^{-1}\cup D_2S_2s,$$ since $s$ commutes with every element of $S_1$ and is inverted by every element of $S_2$. Since $G_1 = C_G(s)$ is a subgroup of $G$ that contains $s$, we have that $G_1 = G_1s$. Similarly, $s^{-1}\in G_1$ and $G_2$ is either empty or the unique nontrivial coset of $G_1$ in $G$, so $G_2s^{-1}=G_2$. Grouping the terms above, we have 
    $$(D_1s\cup D_2s^{-1})S = (D_1S_1\cup D_2S_2)s \cup (D_1S_2\cup D_2S_1)s^{-1}$$
    $$ = G_1s\cup G_2s^{-1} = G_1\cup G_2 = G.$$
    Hence the defenders do occupy a dominating set after moving.
    
Now suppose $i=2$. Then after moving, the defenders occupy the vertices $D_1s^{-1} \cup D_2s$. 
Then $$(D_1s^{-1} \cup D_2 s)S  = D_1s^{-1}S_1 \cup D_1s^{-1}S_2 \cup D_2sS_1 \cup D_2sS_2$$
    $$ = D_1S_1s^{-1} \cup D_1S_2s \cup D_2S_1s\cup D_2S_2s^{-1},$$ since $s$ commutes with every element of $S_1$ and is inverted by every element of $S_2$. Since $G_1 = C_G(s)$ is a subgroup of $G$ that contains $s^{-1}$, we have that $G_1 = G_1s^{-1}$. Similarly, $s\in G_1$ and $G_2$ is either empty or the unique nontrivial coset of $G_1$ in $G$, so $G_2s=G_2$. Grouping the terms above, we have 
    $$(D_1s^{-1}\cup D_2s)S = (D_1S_1\cup D_2S_2)s^{-1} \cup (D_1S_2\cup D_2S_1)s $$ $$= G_1s^{-1}\cup G_2s = G_1\cup G_2 = G.$$
    Hence the defenders do occupy a dominating set after moving.
   \end{proof}

\end{document}
